\chapter{Macchine di Turing multi-nastro}

\section{Introduzione}
Le macchine di Turing multi-nastro presentano due o più nastri. Ogni nastro ha

\begin{itemize}
	\item Una propria sequenza di caselle
	\item Una propria testina meccanica
	\item Un proprio simbolo scandito
\end{itemize}
%
Le testine si muovono in maniera indipendente, ma sono comandate da un unico sistema di controllo. A ogni passo, la macchina esegue simultaneamente le operazioni su tutti i nastri.

\subsection{Istruzioni}

Le istruzioni in una macchina di Turing con $k$ nastri sono rappresentate nel seguente modo:

$$
(q, (t_1, \dots, t_k), (o_1, \dots, o_k), p)
$$
%
Quindi, un'azione su una macchina multi-nastro, consiste nell'esecuzione in parallelo sui $k$ nastri delle operazioni $(o_1, \dots, o_k)$ e nella transizione $q\to p$.

\subsection{Configurazione iniziale}
Per un dato input

$$
Z = s_{i_1} \dots s_{i_n}
$$
%
la configurazione iniziale è definita cosi:

\begin{itemize}
	\item Sistema di controllo nello stato iniziale $q_0$
	\item Input $Z$ contenuto \textbf{nel primo nastro}
	\item La testina del primo nastro sulla prima casella dell'input
	\item Tutti gli altri nastri contengono solo $\square$.
\end{itemize}
%
All'arresto della macchina, l'output viene estratto dal primo nastro. \it Ovviamente queste sono tutte convenzioni, non ci sarebbero implicazioni sulla potenza di calcolo se questi dettagli cambiassero.\rm

\subsection{\textit{k}-istruzioni}
$k>0$.
Una $k$-istruzione è una $k$-upla


\newpage
\section{Simulazione con un solo nastro}
Tutto ciò che fa una macchina di Turing a più nastri è simulabile con un solo nastro (in un trade-off nel numero di azioni da eseguire).

\paragraph{L'idea}
Rappresentare su un singolo nastro:
\begin{itemize}
	\item il contenuto di tutti i nastri, ponendoli in sequenza in maniera opportuna;
	\item la posizione di ciascuna testina;
	\item lo stato corrente della macchina simulata.
\end{itemize}
%
Questo intrudce costi aggiuntivi in termini temporali, ma non cambia la classe delle funzioni calcolabili.

\subsection{Simboli aggiuntivi per la simulazione}

Indicheremo con $\mathfrak{M}$ una $k$-Macchina di Turing, con $\mathfrak{T}$ la macchina simulante, con $t^\rhd$ i simboli detti \textbf{head-marker}, un separatore \#, e...


\subsection{Stati della macchina simulante}

$\mathfrak{T}$ usa due tipi di stati.

\paragraph{Stati semplici} Usati per fasi ausiliarie della simulazione, come inizializzazione e terminazione della codifica.

\paragraph{Stati composti} Servono per la simulazione vera e propria di un passo di $\mathfrak{M}$.
Useremo quindi stati nella forma $(q, U)$ e $(p, V, W)$. Questi sono i \textbf{nomi} degli stati, non coppie o terne.

\begin{itemize}
	\item Gli stati del tipo $(q, U)$: $q$ è un simbolo di $\mathfrak{M}$, mentre $U$ memorizza i simboli già letti dalle testine durante la scansione da sinistra verso destra.
	
	\item Gli stati del tipo $(p, V, W)$: $p$ è il nuovo stato da assumere, $V$ memorizza le operazioni ancora da simulare, e $W$ è una stringa di controllo usata per simulare la singola operazione corrente.
\end{itemize}

Ad esempio, lo stato $(q,\varepsilon)$ indica che $\mathfrak{T}$ sta simulando lo stato $q$ di $\mathfrak{M}$, e non ha ancora letto alcun head-marker.

\subsection{Codifica di una \textit{k}-configurazione}
Sia 

$$
\gamma = ...
$$
%
una $k$-configurazione di $\mathfrak{M}$. La configurazione corrispondente di $\mathfrak{T}$ è data da